\documentclass[11pt]{article}
\usepackage{graphicx} 
\setlength{\parindent}{0pt}
\usepackage{hyperref}
\usepackage{enumitem}
\usepackage[utf8]{inputenc} 
\usepackage[T1]{fontenc}
\usepackage[spanish]{babel}
\usepackage{lipsum}
\usepackage[left=1.06cm,top=1.7cm,right=1.06cm,bottom=0.49cm]{geometry}
\hypersetup{
    hidelinks
}

\begin{document}
\begin{minipage}[c]{0.75\textwidth}
    {\huge \textbf{Ignacio Gomez Barrios}} \\[8pt]
    Córdoba, Argentina 5000 \\
    ignas.gomezb12@gmail.com \textbullet \ +54 9 351 732-5350 \\
    lettu3.github.io
\end{minipage}
\begin{minipage}[c]{0.25\textwidth}
    \begin{flushright}
        \includegraphics[width=3cm]{foto.jpg} 
        \end{flushright}
\end{minipage}

\vspace{5pt}
\hrulefill
\vspace{10pt}

\begin{center}
    \textbf{Perfil Profesional} 
\end{center}
Desarrollador Full Stack (Java y TypeScript) con experiencia en arquitectura backend y desarrollo de interfaces. Combino resolución técnica con comunicación clara tanto con equipos de trabajo como con clientes. Continúo mi formación académica en Ciencias de la Computación en FaMAF.

\vspace{0.5pt}

\begin{center}
    \textbf{Formación} 
\end{center}
\textbf{Universidad Nacional de Córdoba}  \hfill Córdoba, Argentina 

Lic. en Ciencias de la Computación, FAMAF.  \hfill 2022 – En curso 

Materias relevantes: Algoritmos y Estructuras de Datos, Ingeniería del Software 

\vspace{12pt}

\textbf{Instituto Privado de Enseñanza San Agustin} \hfill Córdoba, Argentina 

Bachiller en Economía y Administración  \hfill Graduación: 2020 

\vspace{12pt}

\begin{center}
    \textbf{Experiencia} 
\end{center}
\href{https://inverte.com.ar/}{\textbf{Inverte}}  \hfill Córdoba, Argentina 

\textbf{FullStack Engineer}  \hfill Diciembre 2025 – Julio 2026 
\begin{itemize}[noitemsep, topsep=0pt, partopsep=0pt, parsep=0pt]
    \item Lidero la arquitectura y el desarrollo Backend de una plataforma integral de gestión de stock de inmuebles, utilizando Java Spring Boot, JPA y PostgreSQL. 
    \item Expando las funcionalidades \textit{core} del sistema mediante el desarrollo de hitos clave: un módulo CRM y una nueva suite de Inteligencia Artificial para la generación automatizada de contenido orientado al \textit{Real Estate}. 
    \item Coordino la integración estructural de estos servicios con módulos complementarios del ecosistema (como analíticas de Google), garantizando la escalabilidad y eficiencia general de la plataforma.
\end{itemize}
\vspace{4pt}
\textit{Referencia:} Santiago Siciliano (CEO) -- Tel: +54 9 3517039723

\vspace{12pt}

\href{https://inverte.com.ar/}{\textbf{Inverte}}  \hfill Córdoba, Argentina 

\textbf{Frontend Developer}  \hfill Abril 2025  – Diciembre 2025 
\begin{itemize}[noitemsep, topsep=0pt, partopsep=0pt, parsep=0pt]
    \item Desarrollo de interfaces escalables con React y TypeScript. 
    \item Optimización de carga y rendering de modelos 3D para navegadores. 
    \item Integración de herramientas de gestión comercial con motores de render en tiempo real. 
\end{itemize}

\vspace{12pt}

\begin{center}
    \textbf{Stack} 
\end{center}

\textbf{Lenguajes y frameworks:} Java (Spring Boot, JPA), JavaScript/TypeScript (React, Redux, Vite, Node.js), Python  (FastAPI, Flask), Haskell, C/C++. 

\textbf{Bases de datos:} PostgreSQL, MySQL, SQLite. 

\textbf{Herramientas:} Git, Docker, testing en Java y JavaScript/React, Vite, Webpack, Office (Word, Excel, PowerPoint), Google Workspace (Docs, Sheets, Slides). 

\textbf{Idiomas:} Inglés (técnico), Español (Nativo), Japonés (básico N6 Tokimeki Level Test). 

\newpage

\begin{center}
    \textbf{Certificaciones} 
\end{center}

\href{https://verify.skilljar.com/c/58u875uhh4rp}{\textbf{Claude 101}} \hfill Anthropic, Abril 2026 \\
Fundamentos del uso de Claude para la optimización de tareas diarias, comprensión de sus características principales y bases para la exploración de capacidades avanzadas.

\vspace{12pt}

\href{https://www.credly.com/badges/76745fcb-3f49-4047-a15a-ed79393a3a26/}{\textbf{Introduction to Cybersecurity}} \hfill Cisco, Febrero 2024 \\
Conocimientos fundamentales en ciberseguridad, abarcando el análisis de vulnerabilidades, la detección de amenazas y las estrategias de defensa.

\vspace{12pt}

\href{}{\textbf{N6 Tokimeki Level Test}} \hfill Centro de Cultura e Idioma Japonés en la Argentina, 2025 \\
Certificación de nivel básico de idioma Japonés.

\vspace{12pt}


\end{document}
